\begin{lemma}[Front member uniqueness]
Let $F$ be a family of finite subsets which is a front on $X$, let $x$ be
an increasing enumeration, and let $s,t$ be finite subsets such that
$\mathsf{FrontPrefix}(F,s,x)$ and
$\mathsf{FrontPrefix}(F,t,x)$.
Then $s=t$.
\end{lemma}
\begin{proof}
By \noderef{FrontPrefix}, the hypotheses give $s\in F$, $t\in F$, and
\(\mathsf{ProperPrefixSet}(s,\operatorname{range}(x))\) and
\(\mathsf{ProperPrefixSet}(t,\operatorname{range}(x))\).  Unpacking these
two proper-prefix assertions using \noderef{ProperPrefixSet}, choose
$a,b\in\operatorname{range}(x)$ such that
\[
s=\{k\in\operatorname{range}(x)\mid k<a\},\qquad
t=\{k\in\operatorname{range}(x)\mid k<b\}.
\]
We compare the two cutoffs.

If $a=b$, the two displayed descriptions have the same right-hand side,
so $s=t$.  If $a<b$, then $a\in\operatorname{range}(x)$ and $a<b$, hence
$a\in t$.  For every natural number $k$,
\[
k\in t\text{ and }k<a
\quad\Longleftrightarrow\quad
k\in\operatorname{range}(x)\text{ and }k<b\text{ and }k<a
\quad\Longleftrightarrow\quad
k\in\operatorname{range}(x)\text{ and }k<a
\quad\Longleftrightarrow\quad k\in s.
\]
Thus $s=\{k\in t\mid k<a\}$, and since $a\in t$ this is precisely
$\mathsf{InitialSegment}(s,t)$ by \noderef{InitialSegment}.  Finally, if
$b<a$, the symmetric argument gives $b\in s$ and
$t=\{k\in s\mid k<b\}$, so
$\mathsf{InitialSegment}(t,s)$ by \noderef{InitialSegment}.

In the equality case, \(\mathsf{InitialSegment}(s,t)\) also holds by the
equality clause of \noderef{InitialSegment}.  Therefore in all cases one
of $\mathsf{InitialSegment}(s,t)$ or
$\mathsf{InitialSegment}(t,s)$ holds.  Applying the prefix-free clause of
the front hypothesis \noderef{Front} to the already established members
$s,t\in F$ forces $s=t$.
\end{proof}
